% Part: normal-modal-logic
% Chapter: syntax-and-semantics
% Section: relational-models

\documentclass[../../../include/open-logic-section]{subfiles}

\begin{document}

\olfileid[id]{nml}{syn}{rel}

\olsection{Model Relasional}

Konsep dasar semantik bagi logika modal normal ialah \emph{model
relasional}. Model ini terdiri atas suatu himpunan dunia yang dihubungkan oleh
``relasi aksesibilitas'' biner, beserta suatu penetapan yang menentukan
!!{propositional variable} mana yang dianggap ``benar'' pada dunia mana.

\begin{defn}
  Suatu \emph{model} bagi bahasa modal dasar adalah tripel $\mModel{M}
  = \tuple{W, R, V}$, dengan
  \begin{enumerate}
  \item $W$ suatu himpunan ``dunia'' tak kosong,
  \item $R$ suatu relasi aksesibilitas biner pada~$W$, dan
  \item $V$ suatu fungsi yang menetapkan kepada setiap !!{propositional
    variable}~$p$ suatu himpunan dunia mungkin $V(p)$.
  \end{enumerate}
  Jika $Rww'$ berlaku, kita mengatakan bahwa $w'$ \emph{dapat diakses
    dari}~$w$. Jika $w \in V(p)$, kita mengatakan bahwa $p$ \emph{bernilai
    benar pada dunia}~$w$.
\end{defn}

Keunggulan besar semantik relasional ialah bahwa model dapat direpresentasikan
melalui diagram sederhana, seperti diagram dalam \olref{fig:simple}. Dunia
direpresentasikan oleh simpul, dan dunia~$w'$ dapat diakses dari~$w$ tepat
ketika terdapat panah dari~$w$ ke~$w'$. Selain itu, kita melabeli suatu simpul
(dunia) dengan~$\mTrue{p}$ jika $w \in V(p)$, dan jika tidak,
dengan~\mFalse{p}. \olref{fig:simple} merepresentasikan model dengan
$W = \{w_1, w_2, w_3\}$, $R = \{\tuple{w_1, w_2},\tuple{w_1,
  w_3}\}$, $V(p) = \{w_1, w_2\}$, dan $V(q) = \{w_2\}$.

\begin{figure}
  \begin{center}
    \begin{tikzpicture}[modal]
      \node[world] (w1) [label={[align=right]right:\mTrue{p}\\ \mFalse{q}}]{$w_1$}; 
      \node[world] (w2) [label={[align=right]right:\mTrue{p}\\ \mTrue{q}},
        above right=of w1]{$w_2$}; 
      \node[world] (w3) [label={[align=right]right:\mFalse{p}\\ \mFalse{q}},
        below right=of w1] {$w_3$};
      \draw[->] (w1) to (w2);
      \draw[->] (w1) to (w3);
    \end{tikzpicture}
  \end{center}
  \caption{Sebuah model sederhana.}
  \ollabel{fig:simple}
\end{figure}

\end{document}
